\begin{papercard}{rodriguez2024}
\cf{Stages} \stagetag{L} (LLM extraction of privacy practices from policy text; prompt, parameter and model tuning; determinism checks), \stagetag{I} (comparison with human ground truth using accuracy, precision, recall, F1; no chance-corrected agreement), \stagetag{D} (choice of annotated corpora). A 14-page arXiv paper (May 2024) arguing that ChatGPT (GPT-4 Turbo) and Llama 2 can replace trained NLP pipelines for policy analysis. It is mainly a prompt-engineering and benchmarking study, useful as a worked example of how an LLM annotation setup was tuned and evaluated, and as a source of caveats.

\cf{Claim and setting} Traditional symbolic and statistical NLP for privacy policies needs ``labor-intensive and potentially error-prone manually annotated datasets''; the paper ``proposes the application of Large Language Models (LLMs) as an alternative for effectively and efficiently extracting privacy practices'' and reports F1 above 93\,\% with lower cost and less expertise \pc{rodriguez2024}{1}{This research proposes the application of Large Language Models (LLMs) as an alternative for effectively and efficiently extracting privacy practices from privacy policies at scale}. Task: for a given policy, state which personal-data types are said to be collected or shared (a policy-level yes/no per data type, ``at least once in the policy'').

\cf{Method: design science with split testing} The authors follow Design Science Research with an iterative split-testing process: change one element of the prompt, parameter or model, measure, keep what helps, and benchmark the next test against the updated configuration \pc{rodriguez2024}{3}{We rely on the Design Science Research (DSR) methodology [38] to propose a ChatGPT framework for privacy policy analysis and evaluate its effectiveness}. Ground truth: the MAPP corpus (64 Google Play policies annotated by legal experts), stratified into an \emph{experimental} set of 33 policies used for tuning and a \emph{control} set of 31 policies reserved for validation \pc{rodriguez2024}{3}{retaining an experimental set on which to apply changes and measure their impacts and a control set to validate the final configuration’s overall performance}; stratification checked by annotation counts per policy (means 8.13 and 8.18). Other corpora: OPP-115 (115 policies), APP-350 (350 policies), IT100 (100 policies annotated for international transfers) used for comparison and generalisation.

\cf{Prompt-design experiments (Table 1, GPT-4 Turbo, temperature 0)} Baseline prompt = Data (the policy) + Task (the data type to find) + Output-format instruction; accuracy .793, precision .785, recall .922, F1 .848 \pc{rodriguez2024}{4}{This baseline prompt is structured in Data, Task and Output format instruction segments}. Tests were cumulative \pc{rodriguez2024}{4}{each technique that demonstrated a performance improvement was systematically integrated into the subsequent tests}: stating that ``the following text is a privacy policy'' and quoting the policy (F1 .848 $\to$ .863); placing the policy at the end (slightly worse); \emph{augmenting the task with the definitions used in the manual annotation guidelines} (F1 .887, biggest single gain) \pc{rodriguez2024}{4}{the prompt was augmented to include the internal definitions used for manual annotations in the MAPP dataset}; splitting policy and task into two messages (.888); removing non-relevant paragraphs (no gain); \emph{segmentation}: paragraph-by-paragraph (.889, but costlier), one data type per request (.829) and both (.523, recall $-59.6\,\%$), which suggests that asking for each practice individually loses context \pc{rodriguez2024}{5}{asking for each practice individually may lead to a loss of a broader contextual understanding of the model}; \emph{few-shot} with 1, 2 and 3 randomly chosen annotated paragraphs (.891, \textbf{.902}, .889), so two shots were best \pc{rodriguez2024}{5}{The best result was obtained with two-shot examples}. Adding a persona system message (``extensive knowledge in data protection and law'') did \emph{not} help and lowered results \pc{rodriguez2024}{7}{Neither of the two system instructions improved the results obtained but rather worsened them}.

\cf{Parameters and reproducibility} Temperature 0 and top-$p$ 1 chosen; top-$p$ made little difference. Determinism check: the same prompt on all 33 policies, on 3 days and 5 times of day = 495 requests, 52 discrepancies, i.e.\ 89.5\,\% consistent (vs.\ 59.6\,\% at default temperature 1.0) \pc{rodriguez2024}{6}{we observed 52 discrepancies (i.e., different responses compared to the typical answer), which means 89.5\% consistency in ChatGPT responses}; the \texttt{seed} option plus temperature 0 gave 90.51\,\%, seed alone 84.65\,\% \pc{rodriguez2024}{6}{the combination of 0 temperature and seed shows 90.51\%, being the most reliable combination of these two parameters}. At higher temperature outputs often reported only the first data type asked about \pc{rodriguez2024}{6}{responses frequently reported only the initial data type queried}. Later: GPT-3.5 reached 99.19\,\% determinism, higher than GPT-4 Turbo \pc{rodriguez2024}{13}{the determinism observed in ChatGPT-3.5 (99.19\%) significantly exceeds that of GPT-4 Turbo}.

\cf{Fine-tuning} gpt-3.5-turbo-0613 with a 4,096-token limit required chunking the policy: 73 training chunks aligned with MAPP annotations, 3 epochs \pc{rodriguez2024}{7}{This process led to the creation of a training set comprising 73 chunks, aligned with the manual annotations from the MAPP corpus}. Fine-tuning raised F1 from .670 to .803 and accuracy from .677 to .867 but did not beat GPT-4 Turbo on whole policies (F1 .935), so few-shot prompting is preferred.

\cf{Validation} The final configuration (policy in the first message, task with definitions, output format and two-shot examples in the second; GPT-4 Turbo, temperature 0, fixed seed, top-$p$ 1) gives on the unseen MAPP control set accuracy .916, precision .898, recall .976, F1 .935 \pc{rodriguez2024}{7}{This configuration yields an accuracy of 0.916, a recall of 0.976, a precision of 0.898, and an F1 score of 0.935 on the control set of the MAPP corpus}, and on OPP-115 accuracy .904, recall .912, precision .949, F1 .930 \pc{rodriguez2024}{7}{This validation yielded consistent results: 0.904 accuracy, 0.912 recall, 0.949 precision, and 0.930 F1 score}.

\cf{Comparisons} \emph{Llama 2 (70B-Chat)}: MAPP F1 .882, OPP-115 F1 .753; few-shot examples hurt (format violations, overfitting to the examples) and a 4,096-token limit forced chunking; the authors conclude it ``does not generalize well'' \pc{rodriguez2024}{8}{Llama 2 obtains worse performance against this new dataset, suggesting that, unlike ChatGPT, this Llama 2 configuration does not generalize well to different datasets}. \emph{SVC classifiers on APP-350} for ten fine-grained data types: ChatGPT average F1 84.1\,\% vs.\ 86\,\%, excluding the SIM-serial anomaly (ChatGPT recall 13\,\%) \pc{rodriguez2024}{9}{Excluding the analysis of the SIM serial identifier, which was identified as an anomaly, the comparative evaluation yielded an average F1 score of 84.1\% for ChatGPT, as opposed to 86\% achieved by the SVC-based classifiers}; the authors trace the anomaly to human annotators coding ``device serial number'' under ``SIM serial number''. \emph{PolicyLint} (symbolic, ontology-based): higher precision but low recall, so ChatGPT's F1 is clearly higher; two authors manually mapped PolicyLint's data objects to MAPP categories independently, then reconciled \pc{rodriguez2024}{9}{This matching was independently conducted, followed by an agreement phase for resolving discrepancies}; evaluation counts whether a data type is found \emph{at least once in a policy} \pc{rodriguez2024}{10}{our methodology, where we assess whether a data type is identified at least once in a policy, instead of PolicyLint’s validation across all relevant statements}. \emph{Generalisation} to international-transfer disclosures (IT100) vs.\ an SVM: ChatGPT better on most metrics.

\cf{Cost and speed} Annotators: MAPP used three experts per policy at 1 h 52 min each \pc{rodriguez2024}{11}{three experts annotated each policy, averaging 1 hour and 52 minutes each}; with an assumed \$10/h legal-expert rate, 100 training policies cost up to \$5,601; GPT-4 Turbo overtakes the classifier approach in cost only at around 81,500 policies (GPT-3.5 at about 825,000). Throughput: an average policy needs 6,652 tokens, so rate limits allow about 45 (GPT-4 Turbo) or 150 (GPT-3.5) policies per minute \pc{rodriguez2024}{12}{With an average of 6,652 tokens required in average to fully process a privacy policy, GPT-4 Turbo can analyze up to 45 policies per minute}.

\cf{Discussion} Human annotation errors and disagreements between corpora (e.g.\ attribute-level agreement is lower than for top-level categories, Wagner et al.) are said to undermine supervised classifiers \pc{rodriguez2024}{11}{the average agreement among human annotators for attribute values is considerably lower than for top-level categories}; the paper concludes the LLM approach ``potentially eliminates the reliance on annotated datasets'' \pc{rodriguez2024}{14}{This approach potentially eliminates the reliance on annotated datasets}.

\cf{Limits (card)} Evaluation unit is coarse (data type present anywhere in the policy), not span, statement or purpose; only 31 control policies for the headline F1; tuning and test policies come from the same corpus and the same annotators' definitions are pasted into the prompt, so the benchmark rewards matching those definitions; no confidence intervals, no repeated-run variance for F1 (only for determinism); no chance-corrected agreement ($\kappa$/$\alpha$), so F1 against one reference is not comparable with a human--human $\alpha$; the ``ground truth'' is treated as correct even though the paper itself shows annotation errors; results are tied to specific model versions (GPT-4 Turbo, Llama 2), a closed API that changes, and 2024 prices; reduced-cost claim ignores prompt-development effort; does not test hard categories such as purposes, retention, or negation; no demonstration of the policy-to-behaviour link.

\cf{Use for this thesis} (1) A concrete, citable template for stage L: split the annotated material into a tuning set and an untouched control set; log every prompt change with its metric effect; include the codebook definitions in the prompt; test few-shot count; fix temperature and seed; repeat identical requests to measure determinism; save model name and date. (2) Evidence that a persona instruction and over-fine task decomposition can \emph{reduce} accuracy, so the thesis should test rather than assume. (3) The cost--accuracy--latency trade-offs as one input to a feasibility judgement, with the caveat that numbers are from 2024 and that API drift means any LLM run in the thesis must be dated and stored. (4) Warning for stage I/L: report per-code metrics and, where possible, chance-corrected agreement against two humans instead of F1 against a single ground truth; use the paper's SIM-serial example as an illustration that disagreement may be caused by the human reference. (5) The paper supports treating LLM output as a candidate annotation to be checked by humans, not as ground truth.
\end{papercard}
